\section{Casos de uso reales de producción remota}
\label{sec:casos_uso_remi}

La producción remota ha dejado de ser una propuesta experimental para convertirse en el
modelo operativo de referencia en eventos de gran escala. En esta sección se recogen
cinco casos de uso representativos que ilustran el rango de aplicaciones, desde eventos
masivos con presupuesto de televisión pública hasta proyectos europeos de investigación
con participación de la UPM.

\subsection{Chaos Communication Congress (C3VOC)}

El Chaos Communication Congress (CCC)\acused{CCC} es uno de los mayores congresos de tecnología y
seguridad informática del mundo, celebrado anualmente en Alemania con más de 17\,000 asistentes en sus ediciones de
mayor afluencia~\cite{ccc_36c3_2019} y múltiples salas en emisión simultánea. El equipo de producción de vídeo del
congreso, C3VOC, desarrolló Voctomix~\cite{voctomix_github} precisamente para cubrir
esta necesidad: un mezclador \textit{software} de código abierto capaz de gestionar varias salas
en paralelo sin depender de \textit{hardware broadcast} propietario.

Cada sala del congreso dispone de su propio servidor con una instancia de voctocore,
alimentado por cámaras que envían señal RAW~I420 sobre TCP/MKV y operado remotamente
por un realizador a través de voctogui. Las grabaciones finales y la emisión en directo
se publican en el portal de media del congreso (\url{media.ccc.de}). Este caso de
uso es el origen del proyecto: Voctomix~2.0 hereda la
arquitectura cliente-servidor de voctocore originalmente diseñada para este entorno.

\subsection{Juegos Olímpicos de Tokio 2020}

Los Juegos Olímpicos de Tokio 2020, celebrados en 2021, constituyeron uno de los mayores
despliegues de producción remota en la historia del olimpismo. Olympic Broadcasting
Services implementó un modelo REMI a gran escala en el que la señal de las
cámaras de todas las sedes olímpicas se transportó a un centro de producción centralizado
en Tokio, reduciendo drásticamente el número de equipos y personal técnico desplazado
a cada instalación~\cite{obs_tokyo_2021}.

Este despliegue validó la viabilidad técnica del modelo REMI para eventos con decenas de
sedes simultáneas, señales de alta resolución y requisitos de fiabilidad \textit{broadcast}. La
reducción de la huella de carbono asociada al transporte de equipos fue uno de los
argumentos adicionales que aceleraron la adopción del modelo remoto como estándar para
ediciones futuras.

\subsection{UEFA Champions League (Telefónica Servicios Audiovisuales)}

Telefónica Servicios Audiovisuales implementó producción REMI para la retransmisión de
partidos de la UEFA Champions League jugados en el estadio Santiago Bernabéu de
Madrid~\cite{tsa_remi_2022}. Las cámaras del estadio enviaban la señal mediante SRT
sobre Internet pública hasta el centro de producción remoto, donde los realizadores
operaban el sistema como si estuvieran físicamente presentes en el estadio.

Este caso de uso demuestra la madurez del protocolo SRT para contribución \textit{broadcast}
en condiciones de red reales: el cifrado AES-256 garantiza la seguridad del contenido
deportivo durante el transporte, y la corrección ARQ recupera los paquetes perdidos sin
impacto perceptible en la latencia. La operación resultó transparente para el equipo de
realización, que mantuvo su flujo de trabajo habitual sin adaptaciones significativas.

\subsection{Proyecto 5G Media: Producción desde el Edge (Telefónica/RTVE/UPM)}

El proyecto europeo 5G Media, financiado por la Comisión Europea con la participación
de Telefónica, \ac{RTVE} y la Universidad Politécnica de Madrid, desarrolló el servicio
denominado \textit{Producción desde el Edge}~\cite{rtve_5gmedia_2019}. El 4 de
diciembre de 2019 se realizó en la Cineteca del Matadero de Madrid el primer piloto
operativo: tres cámaras conectadas mediante fibra óptica al \textit{edge} de Telefónica
transmitieron la ficción sonora \textit{La Radio es Sueño} de Radio~3, mientras el
equipo de realización operaba desde las instalaciones de \acs{RTVE} en Torrespaña,
suprimiendo el desplazamiento de la unidad móvil de producción.

La latencia de transporte medida en la red de acceso de Telefónica fue de 3--5~ms,
cumpliendo el requisito de tiempo real para producción en directo. Este
resultado demostró que la infraestructura de \textit{edge computing} puede sustituir
funcionalmente a la unidad móvil en eventos de cobertura local. La industria ha documentado
reducciones de coste relevantes en implantaciones REMI de escala comparable~\cite{tmbroadcast_remi,remi_workflow_ibc2023}.

\subsection{Proyectos europeos de investigación de la UPM}

La Universidad Politécnica de Madrid ha liderado una línea de investigación en
producción audiovisual remota a través de sucesivos proyectos europeos que han
evolucionado desde la validación de tecnologías de red hasta la ciberseguridad
aplicada a entornos multimedia distribuidos:

\begin{itemize}
  \item \textbf{5G Media (2017--2020):} proyecto de la Comisión Europea que exploró
        la viabilidad del \textit{edge computing} para la producción audiovisual
        remota. El piloto con RTVE y Telefónica descrito en la sección anterior
        constituyó su demostración operativa más representativa.

  \item \textbf{NEMO (Network and Edge Media Orchestration):} proyecto europeo
        orientado a la orquestación de servicios multimedia en infraestructuras de
        red distribuidas, en cuyo marco se desplegó Voctomix~1.0 en escenarios de
        producción distribuida. La experiencia en arquitecturas de microservicios
        acumulada en NEMO sirvió de antecedente técnico para el posterior desarrollo
        de Voctomix~2.0 en el proyecto CyberNEMO.

  \item \textbf{CyberNEMO:} proyecto europeo centrado en la ciberseguridad aplicada
        a entornos de producción y distribución multimedia. Voctomix~2.0 fue
        desplegado en el clúster Kubernetes del proyecto, donde el \textit{broker} RabbitMQ
        del sistema fue integrado con las colas
        \texttt{vproduction-uc-media} y \texttt{vqprobe-uc-media}, publicando en
        tiempo real eventos de mezclador durante sesiones de producción audiovisual
        técnica. Este despliegue constituyó la primera validación del sistema de
        telemetría de Voctomix~2.0 en un entorno real.
\end{itemize}

La tabla~\ref{tab:casos_uso_remi} resume los cinco casos de uso presentados, indicando
la escala de despliegue, la tecnología clave empleada y la aportación de cada uno al
contexto del presente trabajo.

\begin{table}[H]
  \centering
  \caption{Casos de uso representativos de producción remota audiovisual.}
  \label{tab:casos_uso_remi}
  \small
  \begin{tabularx}{\textwidth}{|l|l|l|X|}
    \hline
    \textbf{Caso de uso} & \textbf{Escala} & \textbf{Tecnología clave} & \textbf{Aportación} \\
    \hline
    CCC / C3VOC          & Múltiples salas   & Voctomix, MKV/TCP    & Origen del proyecto \\
    JJ.OO. Tokio 2020    & Decenas de sedes  & REMI cloud           & Validación olímpica \\
    UEFA CL (Telefónica) & Estadio único     & SRT + AES-256        & REMI sobre Internet \\
    5G Media (UPM/RTVE)  & Red de acceso     & Edge computing, fibra & Latencia de red 3--5~ms \\
    CyberNEMO (UPM)      & Clúster Kubernetes & RabbitMQ, telemetría & Validación académica \\
    \hline
  \end{tabularx}
\end{table}

Los casos anteriores ilustran que el modelo REMI ha alcanzado madurez tecnológica en
entornos de producción real, desde eventos masivos hasta proyectos europeos de
investigación. El diseño e implementación del sistema que aborda este trabajo se
presenta en el Capítulo~\ref{chap:diseno}.
